\section{Introduction}

In this project, we explore the application of reinforcement learning (RL) techniques to the domain of Pokemon battles. 
Our goal is to develop an RL agent capable of learning effective strategies for battling against a heuristic opponent.

The pokemon battle environment is a turn-based game where two players select actions (such as attacking, switching Pokemon, or using items) to defeat their opponent's team.
The action space is also large, as there are multiple moves and strategies to consider at each turn.
This environment presents three main challenges:
\begin{itemize}
    \item \textbf{High-dimensional state and action spaces}: The large number of possible states and actions makes it difficult for the agent to learn effective policies.
    \item \textbf{Partial observability}: The agent does not have access to the full state of the environment, as it cannot see the opponent's Pokemon or their moves.
    \item \textbf{High stochasticity}: The outcome of actions can be highly variable due to factors such as move accuracy, critical hits, and random effects, making it challenging for the agent to learn consistent strategies.
\end{itemize}

Previous approaches range from simple scripted baselines to modern reinforcement learning and search-based agents. 
Common baselines include hand-crafted heuristics such as always picking the highest base-power move or switching when at a clear type disadvantage, as well as purely random policies used for benchmarking. 
Classical methods like tabular Q-learning and softmax exploration have been applied in simplified environments \cite{kalose2018optimal}, showing consistent gains over random agents while remaining limited by state-space size. 
More recent work uses deep RL (DQN, A3C, PPO) on top of Pokémon Showdown via dedicated environments, often finding policy-gradient methods like PPO particularly effective in stochastic, high-dimensional settings. 
Beyond RL, some systems combine minimax-style lookahead with learned value functions or leverage large language models \cite{jain2025large} and offline sequence models trained on battle logs, which excel at capturing long-horizon strategy and opponent modeling but trade off interpretability and computational cost.

We have implemented two PPO-based agents: a simple policy network and a more complex actor-critic architecture.
The report is structured as follows: we first describe the environment \ref{section:environment}., then we detail the models we implemented \ref{section:models}., followed by the training setup \ref{section:training}. and finally we present our experimental results \ref{section:results}.